\section{Introduction}

% \sout{We have another title candidate: Accelerating on-Device LLM Draft Verification Workloads with NPU.}

With the advancements of commodity hardware and algorithms, deploying Large Language Models (LLMs) on mobile devices is becoming increasingly feasible. Many language models tailored for mobile devices have emerged, including Llama 3.2~\cite{llama32}, MiniCPM~\cite{minicpm,minicpm-v}, Gemma~\cite{gemma1,gemma3}. However, these models generally underperform compared to their larger counterparts. A straightforward approach to improve model performance is to scale up the model size, yet this significantly increases memory consumption and bandwidth requirements, posing serious challenges for resource-constrained mobile platforms.

%Given fixed resource budgets, low-bit post-training quantization\cite{gptq,awq,quip,quipsharp} techniques are commonly recognized as useful methods to offset the cost of using larger models, yet recent studies\cite{empirical-qwen3-quant,kumar2024scaling} show that they become less effective as models become more well trained (in other words, having higher knowledge density\cite{densing-law})\ting{I did not see this statement from this 54 paper. There seems some misunderstanding here. Also quant is only one kind of memory optimization method, just mentioning it here is a distraction from your main story. I suggest to remove it. }.

% Recent test-time scaling methods open up new opportunities to strengthen LLM's abilities through more inference-time compute. Many of the latest models\cite{openai_o1,openai_o3,deepseek_r1} now possess the capability to engage in continuous reflection by generating long chain of thoughts during runtime, exhibiting strong reasoning abilities. Another class of well-known runtime computation scaling methods 
Recently, a new paradigm named \textbf{test-time scaling} have introduced new opportunities to enhance LLM capabilities through increased inference-time computation. Parallel test-time scaling methods involve generating multiple paths and selecting the best sample among a number of generation candidates~\cite{cot-sc,wu2024inference,snell2024scaling,huggingface2025testtime,brown2024large,liu2025can1bsurpass405b,tts-survey,chen2025sets}. So far, these methods are limited to cloud or offline settings where computational resources are abundant.

Intuitively, employing test-time scaling techniques to enhance LLM's generation quality on mobile devices may seem impractical. Mobile devices such as smartphones are typically considered resource-constrained, while LLM inference is known for its high resource consumption. On top of this, scaling compute resources at runtime requires even more computation.

However, recent integration of Neural Processing Units (NPUs) in mobile SoCs has begun to shift this landscape. Vendors including Qualcomm, Intel, and AMD have designed and integrated NPUs to accelerate AI workloads~\cite{hexagon,intel-npu,amd-xdna}. These NPUs not only achieve high peak computing power but also undergo rapid evolution: Qualcomm claims that its Hexagon NPU in Snapdragon X Elite delivers 45 TOPS of INT8 performance~\cite{qualcomm_on_device_ai}, while recent generations of AMD NPUs have achieved $3.1\times$\footnote{The value is obtained by dividing the 50 TOPS of the AMD Ryzen AI 9 HX 370 NPU by the 16 TOPS of the AMD Ryzen 7 8845HS NPU.} performance improvements~\cite{amd_processors_specs}. These developments are transforming the computation capabilities of mobile devices.

% However, the reality is that many smartphone SoCs already incorporate Neural Processing Units (NPUs) to accelerate AI workloads\ting{+ "with fast peak performance improvement, e.g., 8gen1 xx TOPs and 8gen3 xx TOPS"}. Companies such as Qualcomm and MediaTek have their own NPU implementations. Qualcomm claims that its Hexagon NPU delivers up to 45 TOPS of INT8 \ting{add citations and also why here and the following both mention the peak performance?} computing power. The presence of NPUs has transformed the landscape of mobile device computation capabilities.

% \ting{Makes the dramatic conflict here "Even the peak performance is fast improve, only xx\% is actually used in decoding, due to the memory intensive feature. Many works still use CPUs for decoding~\cite{}."}

We discover that mobile NPUs achieve high peak performance through dedicated matrix multiplication units that operate on large matrix tiles. However, in typical LLM inference, GEMM operations often degenerate into GEMV during the decoding phase, resulting in low hardware utilization and waste of computing capabilities of the large-tile optimized matrix units. This underutilization presents an opportunity: test-time scaling methods that increase sampling parallelism can leverage this available compute capacity without substantially adding to inference overhead.

Despite this potential, achieving efficient test-time scaling with mobile NPUs faces significant hardware challenges, which we categorize into two aspects:

\textbf{Precision:} Mobile NPUs were originally designed for coarse-grained quantized models and lack native support for fine-grained group quantization, which is essential for modern LLMs deployed in low bits. We observe that models quantized with conventional per-channel methods suffer severe performance degradation on reasoning tasks that are critical for test-time scaling.

\textbf{Efficiency:} While NPUs excel at matrix multiplication, their general-purpose vector units offer limited compute throughput and memory bandwidth. Many key non-matrix computations in LLM inference for test-time scaling must run on vector units, becoming a prominent bottleneck. Furthermore, the mismatch between the wide SIMD vector components and data granularity, coupled with the hardware's memory access limitations, makes it difficult for software to fully utilize the computing power of vector units, further exacerbating the problem. 

% hzx: maybe I should append "thereby advancing the performance-cost Pareto frontier for on-device LLM deployment" to the first sentence. but this might be an overclaim. I try to be conservative here.
To address these challenges, we present an end-to-end LLM inference system that leverages the abundant compute capacity of mobile NPUs to support test-time scaling workloads. To meet on-device resource constraints and precision requirements, we mainly adopt weight-only 4-bit fine-grained group quantization. For the resulting efficiency challenges, our solution incorporates the following key techniques:

% To address these challenges, we present an end-to-end LLM inference system for Hexagon NPUs and test-time scaling workloads that supports fine-grained weight quantization.\ting{rephrase the sentence above. What did you actually realize? "We present an e2e inference system that incorporates NPU's abundant computing resource and LLM's test-time scaling to achieve new accuracy-cost prato frontier for on-device LLM deployment." Your goal is certainly not for "fine-grained quant" By the way, can you add an overall prato frontier figure in intro?} \ting{The following techniques descriptions need to be improved. You need to specify what is the actual new tech.}To tackle precision issues, we adopt a Post-Training Quantization (PTQ) scheme with fine-grained quantization groups. For the resulting efficiency challenges, our solution involves offline and online stages:

\textbf{Hardware-aware Tile Quantization Scheme:} We present a novel matrix and vector unit-aware quantization layout. Through weight layout transformations before and after quantization, we apply fine-grained group quantization on hardware-friendly tiles and align with NPU's memory access patterns, thereby minimizing runtime memory access overhead and maximizing vector compute utilization.

\textbf{Efficient LUT-Based Computation:} We replace complex key operations, including exponential computation in Softmax and the dequantization process in mixed precision GEMM, with efficient table lookup (LUT) instructions, alleviating computation bottleneck on the vector units.

We evaluate our system across three generations of Qualcomm Snapdragon platforms. Our proposed techniques bring up to $19.0\times$ speedup for mixed-precision GEMM and $2.2\times$ acceleration for Softmax compared to baselines, respectively. The results demonstrate the effectiveness of exploiting mobile NPUs for LLM test-time scaling workloads. Notably, we show that test-time scaling achieves state-of-the-art performance-cost trade-offs: using test-time scaling with smaller models can match or even surpass the performance of larger models running without scaling. To the best of our knowledge, this is the first work to explore the feasibility and evaluate the trade-offs of test-time scaling methods for LLMs with NPUs on mobile devices. Our contributions are summarized as follows:

% We evaluated our system on three generations of Qualcomm Snapdragon SoCs, and the results demonstrate the effectiveness of our proposed techniques. Notably, \ting{textbf the following sentences}our approach achieves better\ting{->"SOTA"? Also you need to briefly introduce your result here} accuracy-latency trade-offs using smaller models with test-time scaling than larger models without it. To the best of our knowledge, we are the first to explore running LLM test-time scaling workloads on mobile devices. Our contributions can be summarized as follows.

\begin{itemize}
    \item We analyze the architecture of modern mobile NPUs and identify underutilization of the specialized matrix units during the LLM decoding phases.
    % hzx: The techniques proposed in this paper are not very important, I suppose.
    \item We present two techniques: a hardware-aware tile quantization scheme and LUT-based computations to accelerate LLM test-time scaling on mobile NPUs.
    \item We design and implement an end-to-end LLM inference system\footnote{Our code is available at https://github.com/haozixu/llama.cpp-npu (main repo) and https://github.com/haozixu/htp-ops-lib (op library)} that leverages mobile NPUs to support test-time scaling workloads with minimal dependency on proprietary software stacks.
    \item We demonstrate that test-time scaling can effectively leverage otherwise wasted NPU compute capacity to enhance the generation quality of on-device small language models, achieving Pareto-frontier performance in accuracy and cost compared to traditional model scaling. It opens up new opportunities for deploying LLMs on mobile devices.
    % \item We implement an end-to-end inference system for Hexagon NPUs and test-time scaling workloads with minimal dependencies.\ting{what do you want to emphasize really, definitely not "minimal dependencies" here you should emphasize your techniques}
    % \item We show that it is possible to leverage the NPU's free computation power to enhance on-device SLM's reasoning capabilities via test-time scaling techniques, approaching or even surpassing the performance of larger models in the traditional LLM scaling scheme of using more parameters. \ting{this should be what you achieved "SOTA accuracy-cost prato frontier"}
    % \item We open up new opportunities for deploying small language models on mobile devices. \ting{this should not be part of contributions or combine with the last point}
\end{itemize}
